\section{Scénarios d'exploitation et impacts}\label{sec:exploitation}

Cette section décrit les scénarios autorisés à rejouer lors de la démonstration.
Elle ne conclut à une compromission que si celle-ci est étayée par une trace datée, une capture ou un résultat conservé par le groupe.

\subsection{Injection SQL sur l'application web}\label{subsec:sqli}

Le premier scénario vise une fonctionnalité DVWA recevant une donnée utilisateur et la transmettant à une requête SQL\@.
L'objectif est de déterminer si la validation insuffisante des entrées permet de modifier la logique de la requête et d'obtenir un accès non prévu à des données.
Cette famille de vulnérabilités est notamment décrite dans l'OWASP Top~10~\cite{owasp-top10}.
Deux rejoues complémentaires sont prévues au niveau de sécurité \texttt{Low} : une requête manuelle contrôlée avec \texttt{curl}, puis la détection automatisée avec \texttt{sqlmap}. La première utilise la charge fixe \texttt{1' OR '1'='1} et établit que plusieurs lignes fictives sont retournées. La seconde ne réalise que des vérifications booléennes ; elle ne doit inclure ni énumération ni extraction de données, notamment pas l'option \texttt{--dump}.
Le protocole expérimental devra préciser la page ciblée, le niveau de sécurité DVWA, l'outil employé et la preuve obtenue. La comparaison distinguera le nombre et la diversité des requêtes produites, le résultat de détection et l'alerte éventuelle associée.
L'analyse des impacts distinguera les données effectivement lues ou modifiées, les privilèges acquis et la portée réelle de l'accès.
Une réponse HTTP inhabituelle ne suffit pas à démontrer une exfiltration.

\subsection{Force brute sur l'authentification DVWA}\label{subsec:bruteforce}

Le second scénario porte sur le mécanisme d'authentification de DVWA\@.
Il utilisera uniquement des identifiants de démonstration créés pour la maquette et une liste de taille volontairement réduite.
L'objectif est d'évaluer la résistance aux tentatives répétées et les traces qu'elles génèrent, non de mener une attaque de volume.
Avant l'essai, le groupe définira un critère de succès mesurable : connexion avec un compte de test, nombre de tentatives nécessaire, présence d'une alerte ou absence de mécanisme de limitation.
Les conséquences seront évaluées à partir des droits du compte de test et non supposées.

\begin{table}[H]
    \centering
    \caption{Grille d'analyse des conséquences.}\label{tab:impacts}
    \begin{tabular}{@{}p{0.28\textwidth}p{0.28\textwidth}p{0.31\textwidth}@{}}
        \toprule
        Critère & Injection SQL & Force brute DVWA \\
        \midrule
        Accès obtenu & À compléter avec preuve & À compléter avec preuve \\
        Données concernées & À compléter avec preuve & À compléter avec preuve \\
        Perturbation du service & À mesurer & À mesurer \\
        Trace conservée & Requête, réponse, journal & Horodatage, tentatives, journal \\
        \bottomrule
    \end{tabular}
\end{table}
